% =========================================================================
% English translation (generated by AI using Claude Opus 4.8 (Max effort)).
% This file was translated from Hungarian to English by AI.
% DISCLAIMER: Please review against the original Hungarian .tex files, before relying on it!
% SOURCE: 5. Valószínűségszámítási és statisztikai alapok.tex
% =========================================================================
\documentclass[margin=0px]{article}

\usepackage{listings}
\usepackage[utf8]{inputenc}
\usepackage{graphicx}
\usepackage{float}
\usepackage[a4paper, margin=0.7in]{geometry}
\usepackage{amsthm}
\usepackage{amssymb}
\usepackage{fancyhdr}
\usepackage{setspace}

\onehalfspacing

\newenvironment{tetel}[1]{\paragraph{#1 \\}}{}

\pagestyle{fancy}
\lhead{\it{CS BSc Final Exam Topics}}
\rhead{5. Basics of probability and statistics}

\title{\textbf{{\Large ELTE IK - Computer Science BSc} \vspace{0.2cm} \\ {\huge Final Exam Topics}} \vspace{0.3cm} \\ 5. Basics of probability and statistics}
\author{}
\date{}

\begin{document}
\maketitle

\begin{center}
\fbox{\parbox{0.92\textwidth}{\small \textit{Note: This document was translated from Hungarian to English by AI. Please verify the information against the original Hungarian source before relying on it.}}}
\end{center}

\begin{tetel}{Basics of probability and statistics}
    Discrete and continuous random variables, law of large numbers, central limit theorem. Statistical estimation, classical statistical tests.
\end{tetel}

\section{Kolmogorov probability space}

The triple $(\Omega, \mathcal{A}, P)$, where:
\begin{itemize}
    \item [$\Omega$] nonempty set, sample space, $\omega$ elementary event.
    \item [$\mathcal{A}$] a system of subsets of $\Omega$, $\mathcal{A} \subset 2^{\Omega}, A \in \mathcal{A}$ events, $\mathcal{A}$ a $\sigma$-algebra. \\
          $\sigma$-algebra:
          \begin{enumerate}
              \item $\Omega \in \mathcal{A}$
              \item $A \in \mathcal{A} \Rightarrow \overline{A} = \Omega \backslash A \in \mathcal{A}$
              \item $A_{1}, A_{2}... \in \mathcal{A} \Rightarrow \bigcup_{i=1}^{\infty}{A_{i}} \in \mathcal{A}$
          \end{enumerate}
    \item [$P$] $: \mathcal{A} \to [0,1]$ set function, probability, for which $P(\Omega) = 1, P(A) \geq 0, \forall A \in \mathcal{A}$, and for pairwise exclusive ($A_{i} \cdot A_{j} = \emptyset, i \neq j$) events $A_{1}, A_{2}... \in \mathcal{A}$ we have $P(\bigcup_{i=1}^{\infty}{A_{i}}) = \sum_{i=1}^{\infty}{P(A_{i})}$.
\end{itemize}

\section{Discrete and continuous random variables}

\begin{itemize}
    \item \textit{Random variable}: $\xi : \Omega \to \mathbb{R}$ a measurable function, that is, one for which $\{\omega : \xi(\omega) < x\} \in \mathcal{A}, \forall x \in \mathbb{R}$, where $\mathcal{A}$ is a system of subsets of the sample space ($\Omega$). ($\omega \in \Omega$ elementary event).
    \item \textit{Distribution / distribution function of a random variable}: $F_{\xi}(x) = P(\xi < x)$, $\forall x \in \mathbb{R}$. \\
          Its properties:
          \begin{enumerate}
              \item $0 \leq F_{\xi}(x) \leq 1$
              \item monotone increasing
              \item left-continuous
              \item $\lim\limits_{x \to -\infty}{F_{\xi}(x)} = 0$, $\lim\limits_{x \to \infty}{F_{\xi}(x)} = 1$
          \end{enumerate}
\end{itemize}

\subsection{Discrete random variables}

Its range is at most countably infinite, that is, it consists of elements $\{x_1 ... x_n ... \}$. Then its distribution: $p_k := P(\xi = x_k)$.

\begin{tabular}{|p{2.5cm}|p{4cm}|p{4cm}|c|c|}
    \hline \textbf{Name}                              & \textbf{Interpretation}                                                     & \textbf{Distribution}                                                                      & \textbf{$EX$}   & \textbf{$D^{2}X$}                                     \\
    \hline indicator \newline $Ind(p)$               & Whether an event of probability $p$ occurs or not.                  & $P(X=1) = p$ \newline $P(X=0) = 1-p$                                                   & $p$             & $p(1-p)$                                              \\
    \hline geometric (Pascal) \newline $Geo(p)$     & On which trial an event of probability $p$ first occurs.         & $P(X=k) = p(1-p)^{k-1}$ \newline $k=1,2...$                                            & $\frac{1}{p}$   & $\frac{1-p}{p^2}$                                     \\
    \hline hypergeometric \newline $Hipgeo(N,M,n)$  & Sampling without replacement.                                         & $P(X=k) = \frac{{M \choose k}{N-M \choose n-k}}{{N \choose n}}$ \newline $k=0,1,...,n$ & $n \frac{M}{N}$ & $n \frac{M}{N}(1 - \frac{M}{N})(1 - \frac{n-1}{N-1})$ \\
    \hline binomial \newline $Bin(n,p)$            & Sampling with replacement.                                               & $P(X=k) = {n \choose k}p^{k}(1-p)^{n-k}$ \newline $k=0,1,...,n$                        & $np$            & $np(1-p)$                                             \\
    \hline negative binomial \newline $Negbin(n,p)$ & On which trial an event of probability $p$ occurs for the $n$-th time. & $P(X=k) = {k-1 \choose n-1}p^{n}(1-p)^{k-n}$ \newline $k=n,n+1,...$                    & $\frac{n}{p}$   & $\frac{n(1-p)}{p^{2}}$                                \\
    \hline Poisson \newline $Poi(\lambda)$           & Rare event.                                                          & $P(X=k) = \frac{\lambda^k}{k!}e^{-\lambda}$                                            & $\lambda$       & $\lambda$                                             \\
    \hline
\end{tabular}

\subsection{Continuous random variables}

A random variable $\xi$ is absolutely continuous if there exists a function $f(x)$ for which $F(x) = \int_{-\infty}^{x}{f(t) dt}$. In this case $f(x)$ is a \textit{density function}. ($F(x)$ is the distribution function.) \\
Another formulation: for $\forall a < b$ we have $P(a < \xi < b) = \int_{a}^{b}{f(t)dt}$, $F_{\xi}(x) = P(\xi < x) = \lim\limits_{a \to -\infty}{P(a < \xi \ b)} = \int_{-\infty}^{x}{f(t)dt}$. \\
Properties of the density function:
\begin{enumerate}
    \item $f(x) = F'(x)$
    \item $f(x) \geq 0$
    \item $\int_{-\infty}^{\infty}{f(x)dx} = 1$
\end{enumerate}

\begin{tabular}{|p{3cm}|c|c|c|c|}
    \hline \textbf{Name}                          & \textbf{Distribution function}                  & \textbf{Density function}                                                        & \textbf{$EX$} & \textbf{$D^{2}X$} \\
    \hline uniform \newline $E(a,b)$
                                                 & $\left\{\begin{array} {lr}
            0               & x \leq a     \\
            \frac{x-a}{b-a} & a < x \leq b \\
            1               & b < x
        \end{array}\right.$
                                                 & $\left\{\begin{array} {lr}
            \frac{1}{b-a} & a < x \leq b \\
            0             & otherwise
        \end{array}\right.$
                                                 & $\frac{a+b}{2}$
                                                 & $\frac{(b-a)^2}{12}$                                                                                                                                             \\
    \hline exponential \newline $Exp(\lambda)$
                                                 & $\left\{\begin{array} {lr}
            1 - e^{-\lambda x} & x \geq 0  \\
            0                  & otherwise
        \end{array}\right.$
                                                 & $\left\{\begin{array} {lr}
            \lambda \cdot e^{-\lambda x} & x \geq 0  \\
            0                            & otherwise
        \end{array}\right.$
                                                 & $\frac{1}{\lambda}$
                                                 & $\frac{1}{\lambda^{2}}$                                                                                                                                          \\
    \hline normal \newline $N(m,\sigma^2)$     & $...$                                      & $\frac{1}{\sqrt{2 \pi}\sigma}e^{-\frac{(x-m)^2}{2\sigma^2}}$ $x \in \mathbb{R}$ & $m$           & $\sigma^2$        \\
    \hline standard normal \newline $N(0,1^2)$ & $\Phi(x)=...$                              & $\frac{1}{\sqrt{2 \pi}}e^{-\frac{x^2}{2}}$ $x \in \mathbb{R}$                   & $0$           & $1$               \\
    \hline gamma \newline $\Gamma(\alpha,\lambda)$
                                                 & $...$
                                                 & $\left\{\begin{array} {lr}
            \frac{1}{\Gamma(\alpha)}\lambda^{\alpha}x^{\alpha-1}e^{-\lambda x} & x \geq 0  \\
            0                                                                  & otherwise
        \end{array}\right.$
                                                 & $\frac{\alpha}{\lambda}$
                                                 & $\frac{\alpha}{\lambda^2}$                                                                                                                                       \\
    \hline
\end{tabular}

\subsection{Concepts}

\begin{itemize}
    \item \textit{Convolution}: $X,Y$ independent random variables, their convolution is the random variable $X+Y$.
    \item \textit{Independence}: $P(\xi_1<x_1, ..., \xi_n<x_n) = \prod_{i=1}^{n}{P(\xi_i<x_i)}$ or in the discrete case: $P(\xi_1=x_1, ..., \xi_n=x_n) = \prod_{i=1}^{n}{P(\xi_i=x_i)}$
    \item \textit{Expected value}: $(\Omega, \mathcal{A}, P)$ probability space, $X:\Omega \to \mathbb{R}$ random variable, $EX = \int_{\Omega}{XdP}$, if it exists. In the discrete case $EX = \sum_{k}{x_k \cdot p_k}$, if absolutely convergent. In the absolutely continuous case $EX = \int_{-\infty}^{\infty}{x \cdot f(x)dx}$, if absolutely continuous.
    \item \textit{Variance}: $D^{2}X = E((X-EX)^2) = EX^2-E^{2}X$
    \item \textit{$l$-th moment}: $EX^l = \int_{\Omega}{x^{l}dP}$, if it exists.
    \item \textit{Standard deviation}: $DX = \sqrt{D^{2}X}$
    \item \textit{Covariance}: $cov(X,Y) = E((X-EX)(Y-EY) = E(XY) - EX \cdot EY$. If $cov(X,Y) = 0$, then $X$ and $Y$ are uncorrelated. (Remark: if two random variables are independent, then $cov(X,Y) = 0$, that is, they are uncorrelated; also $cov(X,X) = D^{2}X$.)
    \item \textit{Correlation}: $R(X,Y) = \frac{cov(X,Y)}{DX \cdot DY}$, measures the linear relationship of two random variables. $R>0 \to$ positive, $R<0 \to$ negative; $R^2 \sim 1 \to$ strong, $R^2 \sim 0.5 \to$ medium, $R^2 \sim 0 \to$ weak.
\end{itemize}

\section{Law of large numbers}

\subsection{Weak law}

$X_1, X_2, ...$ independent, identically distributed, $EX_i = m < \infty$, $D^{2}X_i = \sigma^2 < \infty$. \\
$P(\frac{X_1 + ... + X_n}{n} - m \geq \varepsilon) \rightarrow 0$ $(n \to \infty)$ for $\forall \varepsilon > 0$ (stochastic convergence).

\subsection{Strong law}

$X_1, X_2, ...$ independent, identically distributed, $EX_1 = m < \infty$, $D^{2}X_1 = \sigma^2 < \infty$. \\
$\frac{X_1 + ... + X_n}{n} \rightarrow m$ $(n \to \infty)$ with probability 1. \\
Remark: We prove it with the Chebyshev inequality. ($\frac{\sigma^2}{n \varepsilon^2} \to 0$ $(n \to \infty)$)

\subsubsection{Chebyshev inequality}

$EX$ finite. \\
Then $P(|X-EX| \geq \lambda) \leq \frac{D^{2}X}{\lambda^2}$ \\
Remark: Proof with the Markov inequality.

\subsubsection{Markov inequality}

$X \geq 0, c > 0$. \\
Then $P(X \geq c) \leq \frac{EX}{c}$

\subsection{Types of convergence}

$\xi_n \to \xi$, that is, $\xi$ is convergent.
\begin{itemize}
    \item \textit{stochastically}: if for $\forall \varepsilon > 0$ $P(|\xi_n - \xi| \geq \varepsilon) \rightarrow 0$ $(n \to \infty)$.
    \item \textit{with probability 1 (almost everywhere)}: if $P(\omega : \xi_n(\omega) \to \xi(\omega)) = 1$.
    \item \textit{in $L^p$}: if $E(|\xi_n - \xi|^p) \rightarrow 0$ $(n \to \infty)$ ($p>0$ fixed).
    \item \textit{in distribution}: if $F_{\xi_n}(x) \rightarrow F_{\xi}(x)$ $(n \to \infty)$ at every continuity point of the latter.
\end{itemize}
\textit{Their relationships}: convergence with probability 1 and in $L^p$ are the strongest; from these follows stochastic convergence, and from this convergence in distribution.

\section{Central limit theorem}

$X_1, X_2, ...$ independent, identically distributed, $EX_1 = m < \infty$, $D^{2}X_1 = \sigma^2 < \infty$. \\
Then $\frac{X_1 + ... + X_n - nm}{\sqrt{n}\sigma} \rightarrow N(0,1)$ $(n \to \infty)$ in distribution, that is, $P(\frac{X_1 + ... + X_n - nm}{\sqrt{n}\sigma} < x) \rightarrow \Phi(x)$ $(n \to \infty)$.

\section{Statistical space}

The triple $(\Omega, \mathcal{A}, \mathcal{P})$, if $\mathcal{P} = \{P_{\vartheta}\}_{\vartheta \in \Theta}$ and $(\Omega, \mathcal{A}, P_{\vartheta})$ is a Kolmogorov probability space for $\forall \vartheta \in \Theta$.

\subsection{Concepts}

\begin{itemize}
    \item \textit{Sample}: $\underline{\xi} = (\xi_1,...,\xi_n): \Omega \to \Xi \in \mathbb{R}^n$. ($\xi_i$ random variable)
    \item \textit{Sample space}: $\Xi$, the set of possible values of the sample, often $\mathbb{R}^n, \mathbb{Z}^n$.
    \item \textit{Sample [realization]}: $\underline{x} = (x_1,...,x_n)$, a concrete observation.
    \item \textit{Statistic}: $T: \Xi \to \mathbb{R}^k$.
    \item \textit{Fundamental theorem of statistics}: (Glivenko--Cantelli theorem) $\xi_1, \xi_2, ...$ independent, identically distributed with distribution function $F$. Then for the empirical distribution function $F_n$ we have $\sup_{-\infty<x<\infty}{|F_n(x) - F(x)|} \to 0$ ($n \to \infty$) with probability 1.
\end{itemize}

\section{Statistical estimation}

$(\Omega, \mathcal{A}, \mathcal{P})$ statistical space, $\vartheta \in \Theta$, $P_{\vartheta}(\xi_1<x_1,...,\xi_n<x_n) = F_{\vartheta}(\underline{x})$ \\
$T(\underline{\xi})$ is an \textit{estimator} of $\vartheta$ if $T: \mathbb{R}^n \to \Theta$. \\
$T(\underline{\xi})$ is an \textit{estimator} of $h(\vartheta)$ if $T: \mathbb{R}^n \to h(\Theta)$. \\
\textit{Unbiasedness}: $T(\underline{\xi})$ is an unbiased estimator of $h(\vartheta)$ if $E_{\vartheta}T(\underline{\xi}) = h(\vartheta)$ for $\forall \vartheta \in \Theta$. \\
\textit{Asymptotically unbiased}: $T(\underline{\xi})$ is asymptotically unbiased for $h(\vartheta)$ if $E_{\vartheta}T(\underline{\xi}) \to h(\vartheta)$ $(n \to \infty)$ for $\forall \vartheta \in \Theta$. \\
The unbiased estimator $T_1$ is \textit{more efficient} than the unbiased estimator $T_2$ if $D_{\vartheta}^{2}T_1 \leq D_{\vartheta}^{2}T_2$ for $\forall \vartheta \in \Theta$. \\
It is efficient if it is more efficient than every other unbiased estimator. If an efficient estimator exists, then it is unique.

\begin{itemize}
    \item \textit{Maximum likelihood estimation}: Likelihood function: $L(\vartheta, \underline{x}) = \left\{\begin{array} {lr}
                  P_{\vartheta}(\underline{\xi} = \underline{x}) & \text{discrete}      \\
                  f_{\vartheta, \underline{\xi}}(\underline{x})  & \text{abs. cont.}
              \end{array}\right.$ \\
          In the independent case: $L(\vartheta, \underline{x}) = \left\{\begin{array} {lr}
                  \prod_{i=1}^{n}{P_{\vartheta}(\xi_i = x_i)} & \text{discrete}      \\
                  \prod_{i=1}^{n}{f_{\vartheta, \xi_i}(x_i)}  & \text{abs. cont.}
              \end{array}\right.$\\
          $\hat{\vartheta}$ is the maximum likelihood estimator of the unknown parameter $\vartheta$ if $L(\hat{\vartheta},\underline{\xi}) = \max_{\vartheta \in \Theta}{L(\vartheta,\underline{\xi})}$.
    \item \textit{Method of moments estimation}: $\vartheta = (\vartheta_1, ..., \vartheta_k)$, $\xi_1, ..., \xi_n$, $l$-th moment: $M_l(\underline{\vartheta}) = E_{\underline{\vartheta}}\xi_{i}^{l}$, empirical $l$-th moment: $\hat{M}_l = \frac{\sum_{i=1}^{n}{\xi_i^l}}{n}$. \\ $\hat{\underline{\vartheta}}$ is the method-of-moments estimator of $\underline{\vartheta}$ if it is a solution of the system of equations $M_l(\underline{\vartheta}) = \hat{M}_l$, $l=1..k$.
\end{itemize}

\section{Hypothesis testing}

We set up a hypothesis, and examine whether it is true. We accept or reject it. It can be parametric or non-parametric; we can examine the equality of expected values, of standard deviations, their value, the independence of complete event systems. With a goodness-of-fit test we can determine whether the random variables have a given distribution function, and with a homogeneity test whether two samples have the same distribution. \\
$H_0$: null hypothesis, $\vartheta \in \Theta_0$; $H_1$: alternative hypothesis, $\vartheta \in \Theta_1$; $\Theta = \Theta_0 \bigcup \Theta_1$. \\
\textit{One- and two-sided tests}: For a two-sided alternative hypothesis we assume non-equality, for a one-sided one some relation. For a two-sided test we examine the absolute value of the test value, whether it is within the acceptance region; then, for example, for the u-test the given error percentage has to be halved for the computation, since the $\Phi$ function is symmetric about the $y$ axis.
\begin{itemize}
    \item \textit{Statistical test}: $\Xi = \Xi_e \bigcup \Xi_k$ (disjoint sets) acceptance and critical region. This decomposition is the statistical test. If the observation is an element of the critical region, then we reject the null hypothesis; if it is not an element, then we accept it. $T(\underline{x}) = \left\{\begin{array} {lr} 1 & x \in \Xi_k \\ 0 & otherwise \end{array}\right.$
    \item \textit{Type I error}: $H_0$ is true, but we reject it. Its probability: $P_{\vartheta}(\underline{\xi} \in \Xi_k)$, $\vartheta \in \Theta_0$.
    \item \textit{Type II error}: $H_0$ is false, but we accept it. Its probability: $P_{\vartheta}(\underline{\xi} \notin \Xi_k)$, $\vartheta \in \Theta_1$. \\
          The goal is to make these errors as small as possible. The two probabilities can be improved at each other's expense, if the number of observations is fixed.
    \item \textit{Size of the test}: $\alpha$ is the size of the test if $P_{\vartheta}(\underline{\xi} \in \Xi_k) \leq \alpha$, $\vartheta \in \Theta_0$. \\
          $\alpha$ is the exact size of the test if $\sup_{\vartheta \in \Theta_0}{P_{\vartheta}(\underline{\xi} \in \Xi_k) = \alpha}$.
\end{itemize}

\section{Classical statistical tests}

\begin{itemize}
    \item \textit{u-test}: We assume that the sample is normally distributed ($\xi_i \sim N(m, \sigma^2)$), $i=1..n$, and that the standard deviation is known.
          \begin{itemize}
              \item \textit{One-sample}: The null hypothesis is whether the expected value equals a concrete value ($m_0$); in other words we examine whether the sample mean does not differ significantly from $m_0$. So $H_0: m = m_0$, and in the two-sided case $H_1: m \neq m_0$, in the one-sided case for example $H_1: m \geq m_0$ or $H_1: m < m_0$. \\
                    The value of the u-test: $u = \sqrt{n}\frac{\overline{\xi} - m_0}{\sigma}$. If the null hypothesis is true, then this is approximately standard normally distributed. \\
                    We examine the hypothesis with error probability $\varepsilon$; for this we need the value $\Phi(u_{1-\varepsilon}) = 1 - \varepsilon$. \\
                    In the two-sided case we reject $H_0$ if $|u| > u_{1-\frac{\varepsilon}{2}}$, and accept it if $|u| \leq u_{1-\frac{\varepsilon}{2}}$. \\
                    In the one-sided case we examine the cases $u > u_{1-\varepsilon}$ (right) and $u < u_{1-\varepsilon}$ (left); in these cases we reject $H_0$.
              \item \textit{Two-sample}: Here the conditions are the following: $\xi_i \sim N(m_1, \sigma_1^2)$, $i=1..n$ and $\eta_j \sim N(m_2, \sigma_2^2)$, $j=1..m$. The standard deviations are likewise known. $H_0: m_1 = m_2$, and $u = \frac{\overline{\xi} - \overline{\eta}}{\sqrt{\frac{\sigma_1^2}{n} + \frac{\sigma_2^2}{m}}}$. $H_1 : m_1 > m_2$ is the upper (right) sided, and $H_1 : m_1 < m_2$ the lower (left) alternative hypothesis.
          \end{itemize}
    \item \textit{t-test}:	In this test the standard deviation is not known, but we assume a normal distribution just as for the u-test. $\xi_i \sim N(m, \sigma^2)$, $i=1..n$.
          \begin{itemize}
              \item \textit{One-sample}: $H_0: m = m_0$. The alternative hypothesis similarly to the u-test. $t = \sqrt{n}\frac{\overline{\xi} - m_0}{\sqrt{\sigma_{*}^2}}$, where $\sigma_{*}^2$ is the corrected empirical variance, which we can compute from the sample. (Remark: instead of $n$ we divide by $n-1$ in the formula.) This value is $t$-distributed under $H_0$, with $n-1$ degrees of freedom. This test is also called the Student's test.
              \item \textit{Two-sample}: $\xi_i \sim N(m_1, \sigma_1^2)$, $i=1..n$ and $\eta_j \sim N(m_2, \sigma_2^2)$, $j=1..m$. In this case too the standard deviation is not known, but we assume that the standard deviations of the two samples are equal. Then $t_{n+m-2} = \sqrt{\frac{nm(n+m-2)}{n+m}}\frac{\overline{\xi} - \overline{\eta}}{\sqrt{\sum{(\xi_i - \overline{\xi})^2} + \sum{(\eta_j - \overline{\eta})^2}}}$. $n+m-2$ is the degrees of freedom of the test.
          \end{itemize}
    \item \textit{F-test}: Usable for two samples. This test is suitable for examining the equality of standard deviations, so here $H_0 : \sigma_1 = \sigma_2$. If the variances of the two samples are equal, then their quotient tends to $1$. $f_{n-1,m-1} = max(\frac{\sigma_1^2}{\sigma_2^2}, \frac{\sigma_2^2}{\sigma_1^2})$. Of the two degrees of freedom the first belongs to the sample in the numerator of $f$ ($-1$), the second to that in the denominator.
    \item \textit{Welch test}: Also called the d-test. Similar to the two-sample t-test, but here the equality of the standard deviations does not have to be assumed. Its degrees of freedom can be computed with a complicated formula.
    \item \textit{Sequential tests}: $V_n = \frac{\prod{f_1(x_i)}}{\prod{f_0(x_i)}} = \frac{L_1(\underline{x})}{L_0(\underline{x})}$. $f_0$ is the density function under the null hypothesis, $f_1$ that under the alternative hypothesis. Given a value $A$ and a value $B$, $A<B$. If $V_n \geq B$, then we reject $H_0$; if $V_n \leq A$, then we accept it; and if $A < V_n < B$, then we take a new sample element. \\
          By \textit{Stein's theorem} $N$ is finite with probability 1. $N = min\{n: V_n \leq A \vee V_n \geq B \}$.
    \item \textit{Quality control}: We look at $n_1$ elements, with thresholds $c_1 < c_2$ and $c_3$. If $X_1 \leq c_1$, then we accept $H_0$; if $X_1 \geq c_2$, then we reject it. If $c_1 < X_1 < c_2$, then we look at $n_2$ elements, and if $X_1 + X_2 \leq c_3$, then we likewise accept $H_0$. The expected sample size measures its efficiency.
    \item \textit{$\chi^2$-test}: $H_0: A_1,...,A_n$ a complete event system. $P(A_i) = p_i, i=1..n$, $\nu_i$ the frequency. If the null hypothesis holds, then $\frac{\nu_i}{n} \sim p_i$. \\
          $\chi^2 = \sum{\frac{(\nu_i - np_i)^2}{np_i}}$. This is $\chi^2$-distributed, with $r-1$ degrees of freedom. $r$ is the number of groups added together. (Remark: if the frequency in some group were too small, then we merge those.) \\
          The $\chi^2$-test can be used for goodness-of-fit, homogeneity and independence testing. (Remark: there is a different formula for each.)
    \item \textit{Other tests}:
          \begin{itemize}
              \item \textit{Kolmogorov--Smirnov test}: whether 2 empirical distribution functions are equal (homogeneity test), or, for 1 sample, whether it equals some distribution function. $D_{m,n} = \max_{x}{|F_n(x) - G_m(x)|}$. $X_i$ with distribution function $F$, $Y_j$ with $G$. $H_0 : F \equiv G$.
              \item \textit{Sign test}: How many times it holds that something is positive.
              \item \textit{Wilcoxon test}: (rank statistic), to test $P(X > Y) = \frac{1}{2}$ we count for how many pairs it holds that $X_i > Y_j$.
          \end{itemize}
\end{itemize}

\end{document}
